\documentclass[main]{subfiles}
\usepackage[sols]{workbook}
\numbering{ws1-11}

\begin{document}

\ifSubfilesClassLoaded{%
  \chaptermark{\chapterone}
}{}

\section{Improper Integrals III} \label{ws1-11}
\begin{framed}
\textbf{Lesson Goals:}
\begin{itemize}
\item You should be able to determine if an improper integral of an unbounded function on a bounded interval converges or diverges, by converting it to a one-sided limit and evaluating. 
\item You should be able to identify all the ways in which an integral is improper, and split it into a sum of limits of proper integrals.
\item You should be able to use multiple strategies (including the Comparison Theorem) to determine if an improper integral converges or diverges. 
\item You should be able to reason abstractly about the convergence behavior of improper integrals.
\end{itemize}
\end{framed}


 
\begin{problemonly}
\begin{framed}
\centerline{\textbf{General Strategy for Evaluating Improper Integrals}}
An integral is \textbf{improper} if
\begin{itemize}
\item the interval of integration is unbounded and/or 
\item the integrand is unbounded somewhere on the interval of integration
\end{itemize}

The basic strategy for dealing with an improper integral is:
\begin{itemize}[topsep=0pt,parsep=0pt,itemsep=0pt]
\item Identify all improprieties.
\item If necessary, split up the integral into a sum of integrals so
that every impropriety is an endpoint of integration and each new
integral has at most one impropriety.
\item Compute each of these new improper integrals as a (one-sided)
limit of definite integrals.
\item If any one of the new integrals diverges, the original improper
integral diverges. If all of the new integrals converge, sum to find
out what the original converges to.

\vspace{8pt}  \textbf{ Caution}: This means that
divergent summands can never ``cancel'' one another!
\end{itemize}
\end{framed}
\end{problemonly}

\begin{enumerate}

\item 
Carry out steps 1 and 2 of the above strategy on the following
integrals.
\begin{enumerate}
  \item 
  \begin{problem}[\vfill]
  $\displaystyle \int_{-2}^2 \frac{x}{x^2 - 1}~dx$.
  \end{problem}

  \begin{solution}
  The integrand has improprieties at $-1$ and 1.  So, let's split like
  this:
  \[ \int_{-2}^2 \frac{x}{x^2 - 1}~dx = \boxed{\int_{-2}^{-1}
  \frac{x}{x^2 - 1}~dx + \int_{-1}^0 \frac{x}{x^2 - 1}~dx +
  \int_{0}^1 \frac{x}{x^2 - 1}~dx + \int_1^2 \frac{x}{x^2 -
  1}~dx}.\] (This is not the only right answer; you could choose a
  number other than 0 between $-1$ and $1$ to split at.)
  \end{solution}

  \item
  \begin{problem}[\vfill]
  $\displaystyle \int_{-\infty}^\infty \sin x~dx$.
  \end{problem}

  \begin{solution}
  This integral has two improprieties: the $-\infty$ and the $\infty$.
  It doesn't really matter where we split it.  We can pick any real
  number $a$ and use $\displaystyle \int_{-\infty}^\infty \sin x~dx =
  \boxed{\int_{-\infty}^a \sin x~dx + \int_a^\infty \sin x~dx}$.
  \end{solution}

\end{enumerate}

\pspace{\newpage}

\item
Determine whether the following integrals converge or diverge.  Explain
your reasoning.  (You may use the above strategy to actually compute the
integral or use the Comparison Theorem to compare to an integral that is
easier to compute.)
\begin{enumerate}
\item 
\begin{problem}[\vfill]
$\displaystyle \int_0^\infty e^{-x^2}~dx$.
\end{problem}

\begin{solution}
We don't know how to find antiderivatives of $e^{-x^2}$ (in fact, as
we've discussed before, the antiderivatives of $e^{-x^2}$ are not
elementary functions, so we really have no hope of writing them down).
Therefore, we had better use comparison.  

In this case, the integral has only one impropriety, at $\infty$.  So,
in deciding whether the integral converges, we should be looking at
what happens to the integrand when $x$ is very large.  The improper
integral involving exponentials that we are familiar with is the
integral of $e^{-x}$.  When
$x$ is large, $e^{-x^2}$ is smaller than $e^{-x}$, so the fact that
$\displaystyle \int_0^\infty e^{-x}~dx$ converges suggests that
$\displaystyle \int_0^\infty e^{-x^2}~dx$ ought to converge as well.
Moreover, the reasoning we've used suggests that we should be
comparing $e^{-x^2}$ with $e^{-x}$. 

We'd like to say that $0 \leq e^{-x^2} \leq e^{-x}$ in order to do the
comparison.  However, this is only true when $x \geq 1$.  (When $x
\geq 1$, $x^2 \geq x$, so $-x^2 \leq -x$, and $e^{-x^2} \leq e^{-x}$.)
So, let's break up the original integral:
\begin{equation}
\int_0^\infty e^{-x^2}~dx = \int_0^1 e^{-x^2}~dx + \int_1^\infty
e^{-x^2}~dx
\end{equation}
The first summand, $\displaystyle \int_0^1 e^{-x^2}~dx$, is a definite
integral (no improprieties at all), so there is no question that it
converges.  In the second integral, we're now looking only at $x \geq
1$, so we can use comparison: we know that $0 \leq e^{-x^2} \leq
e^{-x}$ when $x \geq 1$, and $\displaystyle \int_1^\infty e^{-x}~dx$
converges by Problem Set 9 \#1(b).  Therefore, the Comparison
Theorem says that $\displaystyle \int_1^\infty e^{-x^2}~dx$ converges.
So both summands on the right side of
(1) converge, which means that
$\displaystyle \int_0^\infty e^{-x^2}~dx$ \fbox{converges}.
\end{solution}



\item 

\note{Sometimes students want to think about this as a $p$-integral;
they say things like, ``If we integrate, we'll get an exponent that's
at least 3.''  A good example to give these students is $\displaystyle
\int \frac{1}{x^2}~dx$ vs. $\displaystyle \int \frac{1}{x^2 + 1}~dx$.}

\begin{problem}[\vfill]
$\displaystyle \int_1^\infty \frac{1}{x^4 + 2}~dx$.
\end{problem}

\begin{solution}
This doesn't look like an integral we actually want to evaluate, so
let's try to use comparison instead.  First, let's just try to decide
whether we think the integral converges or diverges.  Here (in {\ideacolor})
is a good thought process to go through.

\ideas{This integral has only one impropriety, at $\infty$.  So, in
deciding whether the integral converges, we should be looking at
what happens to the integrand when $x$ is very large.  When $x$ is
very large, $x^4$ is much much bigger than 2, so $\frac{1}{x^4 + 2}$
is pretty close to just $\frac{1}{x^4}$.  We know that
$\displaystyle \int_1^\infty \frac{1}{x^4}~dx$ converges, since it's
the $p$-integral with $p = 4$, and we know all $p$-integrals with $p
> 1$ converge.  So, it seems like $\displaystyle \int_1^\infty
\frac{1}{x^4 + 2}~dx$ should converge, too.}

To write a good mathematical argument, we should use the Comparison
Theorem.  To get an idea of whether the integral in question
converged, we compared to $\displaystyle \int_1^\infty
\frac{1}{x^4}~dx$, so a good first try is to compare $\frac{1}{x^4 +
2}$ with $\frac{1}{x^4}$.  This works well in this case:

For any $x \geq 1$, $0 \leq \frac{1}{x^4 + 2} \leq \frac{1}{x^4}$.
Since $\displaystyle \int_1^\infty \frac{1}{x^4}~dx$ converges (as
explained already), the Comparison Theorem tells us that
$\displaystyle \int_1^\infty \frac{1}{x^4 + 2}~dx$ also
\fbox{converges}.
\end{solution}



\item 
\begin{problem}[\vfill]
$\displaystyle \int_0^\infty \frac{1}{e^x + x}~dx$.
\end{problem}

\begin{solution}
It's not clear how to find an antiderivative of $\frac{1}{e^x + x}$,
so we should try comparison.  \ideas{Again, the only impropriety this
integral has is at $\infty$, so the thing we need to understand is
how the integrand behaves when $x$ is really large.  In the
denominator, when $x$ is large, $e^x$ is much bigger than $x$, so
the integrand behaves more like $\frac{1}{e^x}$.  We know by
Problem Set 9, \#1(a) that $\displaystyle \int_0^\infty
\frac{1}{e^x}~dx$ converges, which suggests that this integral
should converge as well.}  We can justify this using the Comparison
Theorem:

Since $0 \leq \frac{1}{e^x + x} \leq \frac{1}{e^x}$ for all $x$ and
$\displaystyle \int_0^\infty \frac{1}{e^x}~dx$ converges by
Problem Set 9, \#1, the Comparison Theorem tells us that
$\displaystyle \int_0^\infty \frac{1}{e^x + x}~dx$ also
\fbox{converges}.
\end{solution}


\item 
\begin{problem}[\vfill]
$\displaystyle \int_{-\infty}^\infty \frac{1}{x^3}~dx$.
\end{problem}

\begin{solution}
$\displaystyle \int \frac{1}{x^3}~dx$ is something we can evaluate
fairly easily, so we should be able to evaluate this improper integral
to decide whether it converges.

This integral has three improprieties: $-\infty$, $0$, and $\infty$.
So let's split it up as $\displaystyle \int_{-\infty}^{-1}
\frac{1}{x^3}~dx + \int_{-1}^0 \frac{1}{x^3}~dx + \int_0^1
\frac{1}{x^3}~dx + \int_1^\infty \frac{1}{x^3}~dx$.  We can compute
each by taking a limit of definite integrals.  We know that
$\displaystyle \int_1^\infty \frac{1}{x^3}~dx$ converges, as it's the
$p$-integral with $p = 3$ (and we know $p$-integrals with $p > 1$ all
converge).  Let's look at the piece before that:
\begin{align*}
\int_0^1 \frac{1}{x^3}~dx & = \lim_{a \to 0^+} \int_a^1
\frac{1}{x^3}~dx \\
& = \lim_{a \to 0^+} \left( \left. - \frac{1}{2} x^{-2}
\right|_{a}^1 \right) \\
& = \lim_{a \to 0^+} \left( -\frac{1}{2} + \frac{1}{2 a^2} \right)
\\
& = \infty
\end{align*}
Since this integral diverges, we don't even need to check the other
pieces; we already know that the original integral $\displaystyle
\int_{-\infty}^\infty \frac{1}{x^3}~dx$ \fbox{diverges}.
\end{solution}


\end{enumerate}
\pspace{\newpage}

 \item
		\begin{problem} Suppose that $f(x)$ is a positive, continuous function.  Suppose we know that $\displaystyle \lim_{x\rightarrow \infty} f(x) = \frac{1}{100}$.  
		What can you say about the convergence of the improper integral $\displaystyle \int_1^\infty \frac{f(x)}{x} \,dx$?\, 

 Please \textbf{circle one} choice, and write a brief explanation.\\

      \begin{multicols}{3}
        Definitely Converges\\
\solonly{\fbox{Definitely Diverges}\\}
\probonly{Definitely Diverges\\}
        Not enough information to tell
      \end{multicols}
		\end{problem}

    \probonly{\textbf{Explain:}}
\vfill
\begin{solution}
Your gut feeling should be that for very large $x$-values, the function $f(x)$ is \textit{very close} to $\frac{1}{100}$, and so the function $\frac{f(x)}{x}$ is \textit{very close to} $\frac{1/100}{x} = \frac{1}{100x}$.  Since $\displaystyle \int_1^\infty \frac{1}{100x}$ is a constant multiple of the $p$-integral with $p=1$, it diverges. So, we should expect that our mystery integral $\displaystyle \int_1^\infty \frac{f(x)}{x}\ dx$ diverges too.

Let's show this carefully using the Comparison Theorem:

Since $\displaystyle \lim_{x\rightarrow \infty}f(x)=\frac{1}{100}$ (i.e., $f(x)$ gets infinitely close to 1/100), there is a number $N$ such that for $x\ge N$ we always have, for example, $\displaystyle f(x) > \frac{1}{200}$.  Then for $x\ge N$ we have: 
$$ 0 \leq \frac{1}{200x}  = \frac{1/200}{x} \le \frac{f(x)}{x}$$

The integral $\displaystyle \int_N^\infty \frac{1}{200x} \ dx$ is a constant multiple of the $p$-integral with $p=1$, so we know it diverges. (Again, the fact that we're starting the integral at $x=N$ instead of $x=1$ does not affect this divergence.) Thus, by comparison, the integral $\displaystyle \int_N^{\infty} \frac{f(x)}{x} \ dx$ diverges. Therefore, the original integral $\displaystyle \int_1^{\infty} \frac{f(x)}{x} \ dx$ diverges as well.
\end{solution}

\item \label{spring2003-midterm1-1c}
      \begin{problem}[\vfill]
        Write an improper integral that gives the volume obtained by revolving
        around the $x$-axis the region bounded below by the $x$-axis, on
        the left by the $x=e$ and above by $\displaystyle y =
        \frac{1}{\sqrt{x} \ln x}$.  Please explain in detail how you
        obtain the your improper integral; be sure your slicing methodology is clear.
    
    Then, determine if your improper integral coverges or diverges. Explain your answer completely.
      \end{problem}
      
      \begin{solution}
        You might not know how to graph $y = \frac{1}{\sqrt{x} \ln x}$
        exactly, but you should be able to see that, for $x \geq e$, it is
        positive and decreasing and that it tends to $0$ as $x \to
        \infty$.  So, here is a sketch of the region in question.
        \begin{center}
          \includegraphics[width=1.5in]{images/spring2003-midterm1-1c-region}
        \end{center}
        Let's first find the volume obtained by rotating the region
        between $x = e$ and $x = b$ (where $b$ is a really big number).
        Then we'll take the limit as $b \to \infty$ to get the volume
        obtained by rotating the whole region.  Let's slice the region
        between $x = e$ and $x = b$ vertically into $n$ slices of equal
        width $\Delta x$ (so we are slicing the interval $[e ,b]$), label
        $x_0 = e, \ldots, x_n = b$ as usual, and look at the $k$-th slice:
        \begin{center}
          \includegraphics[width=1.5in]{images/spring2003-midterm1-1c-region-sliced}
        \end{center}
        When we rotate this slice about the $x$-axis, we get
        (approximately) a disk of radius $\frac{1}{\sqrt{x_k} \ln x_k}$
        and thickness $\Delta x$, so
        \[ \textrm{volume of $k$-th slice} \approx \pi \cdot \frac{1}{x_k
          (\ln x_k)^2} \Delta x.\] Summing over all slices and taking the
        limit as the number of slices $n$ tends to $\infty$, the volume
        obtained by rotating the region between $x = e$ and $x = b$ is
        $\displaystyle \pi \int_e^b \frac{1}{x (\ln x)^2}~dx$, and we are
        interested in the limit as $b \to \infty$ of this integral.  To
        evaluate, we'll substitute $u = \ln x$ (so $du = \frac{1}{x}~dx$)
        \begin{align*}
          \pi \int_e^b \frac{1}{x (\ln x)^2}~dx &= \pi \int_1^{\ln b}
          \frac{1}{u^2}~du \\
          &= - \pi \left( u^{-1} \Big|_1^{\ln b} \right) \\
          &= \pi \left(- \frac{1}{\ln b} + 1 \right)
        \end{align*}
        The limit as $b \to \infty$ of this quantity is just
        $\boxed{\pi}$.
      \end{solution}
\pspace{\newpage}
  \item \label{fall2008-midterm1-3}
  \begin{enumerate}
  \item
    \begin{problem}[\vfill]
      Give an example of a function $f(x)$ which is continuous and defined
      on $[1, \infty)$ such that $\displaystyle \lim_{x\to\infty}f(x)=0$
      and $\displaystyle \int_1^\infty f(x)~dx$ converges.
    \end{problem}

    \begin{solution}
      There are infinitely many such functions; $f(x) = \displaystyle
      \frac{1}{x^2}$ is one.  (An even simpler one is $f(x) = 0$.)
    \end{solution}

  \item
    \begin{problem}[\vfill]
      Give an example of a function $f(x)$ which is continuous and defined
      on $[1, \infty)$ such that $\displaystyle \lim_{x\to\infty}f(x)=0$
      and $\displaystyle \int_1^\infty f(x)~dx$ diverges.
    \end{problem}

    \begin{solution}
      There are infinitely many such functions; $f(x) = \displaystyle
      \frac{1}{x}$ is one that you should be quite familiar with.
    \end{solution}

  \item
    \begin{problem}[\vfill]
      Is it possible to find an example of a function $f(x)$ which is
      continuous and defined on $[1, \infty)$ such that $f(x)>0$,
      $\displaystyle \int_1^\infty f(x)~dx$ diverges and $\displaystyle
      \int_1^\infty f(x)^3~dx$ converges?  If so, give an example. If not,
      explain why not.
    \end{problem}

    \begin{solution}
      $f(x) = \displaystyle \frac{1}{x}$ is such a function; we know that
      $\displaystyle \int_1^\infty \frac{1}{x}~dx$ diverges (it's the
      $p$-integral with $p = 1$, and we know $p$-integrals with $p \leq 1$
      all diverge) and that $\displaystyle \int_1^\infty \frac{1}{x^3}~dx$
      converges (it's the $p$-integral with $p = 3$, and we know
      $p$-integrals with $p > 1$ all converge).
    \end{solution}

  \end{enumerate}           
\end{enumerate}

\end{document}